\section{Multi-Root Workspace}
\label{sec:workspace}

\subsection{Problem}

A developer working on a project often needs to navigate \emph{reference
codebases} alongside their main project: a dependency's source, a shared
library, or a related service. codetopo supports this via \emph{workspace
roots}: named, absolute filesystem paths that are all indexed and queryable
within a single server instance.

\subsection{Design}

All roots share a single \texttt{index.sqlite} database with no mode
switching. The \texttt{roots} table (\S\ref{sec:data-model}) maps each root
to an integer ID; the \texttt{files} table records \texttt{root\_id} on every
file row. There is no separate per-root database.

\subsection{ID offset scheme}

Nodes, files, and edges across all roots share a single ID space. codetopo
ensures uniqueness with a multiplicative offset:

\[
  \text{global\_id} = \text{root\_id} \times 10^9 + \text{local\_id}
\]

\texttt{local\_id} is the SQLite autoincrement rowid within the root's space
(values $1$ through $10^9 - 1$). A signed 64-bit integer ($\approx 9.2 \times
10^{18}$) supports around 9 billion root-slots at this scale, far beyond any
practical need. Because all IDs are stored in a single file, no cross-root
join is needed: a node ID uniquely identifies both the root and the local
record. Figure~\ref{fig:id-offset-array} visualizes one root-local slot before
the offset formula lifts it into the global ID space.

\begin{figure}[htbp]
  \centering
  \tritonnext{width=0.7\textwidth}
\begin{triton}
array
  title local_id within one root slot
  cells 1 2 3 4 5
  index
  ptr local_id -> 2
\end{triton}
  \caption{A root-local ID slot before conversion into the global
           \texttt{root\_id * 10\^{}9 + local\_id} namespace.}
  \label{fig:id-offset-array}
\end{figure}

\subsection{Path handling}

All paths stored in the database are \emph{relative to their root}. Tools that
accept filesystem paths call \texttt{resolve\_db\_path()}, which tries three
strategies in order:
\begin{enumerate}
  \item Absolute path match: \texttt{root.path || '/' || file.path = :abs}.
  \item Exact relative match against \texttt{files.path}.
  \item Suffix \texttt{LIKE} match: \texttt{files.path LIKE '\%' || :suffix}.
\end{enumerate}
The suffix match handles editors that pass workspace-relative paths without
specifying the root. This approach is necessary because workspace roots may
live at different absolute paths on different machines.

\subsection{Workspace operations}

\subsubsection{Adding a root}

\texttt{workspace\_add} accepts an absolute path. The flow is:
\begin{enumerate}
  \item Canonicalize the path (resolve symlinks).
  \item Ensure the path is indexed: if no \texttt{roots} row exists, run a
        full index first.
<<<<<<< HEAD
  \item \texttt{ATTACH} the source database (if separate).\footnote{A SQLite
        gotcha: \texttt{ATTACH DATABASE} must be issued \emph{before}
        \texttt{BEGIN TRANSACTION} on the same connection. Attempting to
        attach inside an open transaction raises \texttt{SQLITE\_ERROR}. The
        workspace-add flow therefore attaches the source database first, then
        opens the write transaction for the row copy.}
=======
  \item \texttt{ATTACH} the source database (if separate).
>>>>>>> origin/main
  \item \texttt{INSERT} a row into \texttt{roots} to obtain the new
        \texttt{root\_id}.
  \item Copy all \texttt{files}, \texttt{nodes}, \texttt{edges}, and
        \texttt{refs} rows, re-keying IDs using the offset formula.
  \item Rebuild \texttt{content\_fts} from disk (line by line per file),
        because contentless FTS tables cannot be copied directly.
\end{enumerate}

\subsubsection{Removing a root}

\texttt{workspace\_remove} runs:
\begin{enumerate}
  \item Delete \texttt{content\_fts} rows for the root's files.
  \item \texttt{DELETE FROM roots WHERE id = :root\_id}.
  \item Foreign-key cascades remove dependent rows in \texttt{files},
        \texttt{nodes}, \texttt{edges}, and \texttt{refs}.
\end{enumerate}
The entire operation runs in a single transaction.

\subsubsection{Listing roots}

\texttt{workspace\_list} returns all roots with path, symbol count, file
count, last-indexed timestamp, and indexing status.

\subsection{Indexer guard}

The normal \texttt{codetopo index} command only re-indexes files belonging to
the main project (\texttt{files.root\_id IS NULL} or the root with
\texttt{id = 0}). It does not overwrite data for workspace roots added via
\texttt{workspace\_add}. This prevents accidental resets of externally indexed
roots.

\subsection{Concurrent access}

SQLite WAL mode allows the Query Engine to serve reads while the Indexer is
writing. The Indexer acquires a write lock only during the commit phase of each
file transaction; reads proceed unblocked during parse and extraction.
